% 전원부의 레일과 안전 회로(한 줄 그림). 그림 fig:hw-power-unit
% 근거: hardware/common/power-unit.md 3-4절, docs/decisions/0003-interfaces.md 4.6, 6, 7절.
\begin{tikzpicture}[x=1cm,y=1cm,
    psu/.style={draw,rounded corners=1.5pt,minimum width=24mm,minimum height=11mm,align=center,inner sep=2pt,font=\scriptsize},
    ld/.style={draw=cgray,fill=white,rounded corners=1.5pt,align=left,inner sep=2.5pt,font=\tiny},
    km/.style={draw=cred,fill=credbg,minimum width=13mm,minimum height=6mm,inner sep=1pt,font=\tiny,align=center},
    sr/.style={draw=cred,fill=white,minimum width=11mm,minimum height=7mm,inner sep=1pt,font=\scriptsize,align=center},
    rail/.style={line width=1.4pt},
    lbl/.style={font=\tiny,inner sep=1pt}]
  % 상용 전원
  \node[psu,fill=clight,draw=cgray,minimum height=16mm] (ac) at (0,0) {상용 전원\\인입구(접지)\\주 스위치\\차단기 10 A};
  % 전원 장치
  \node[psu,fill=cbluebg,draw=cblue] (pl) at (3.0,2.6) {PSU-L\\24.5 V 10 A};
  \node[psu,fill=corangebg,draw=corange] (ps) at (3.0,0) {PSU-SW\\24.5 V 27 A\\{\tiny 릴레이식, 6유닛마다}};
  \node[psu,fill=credbg,draw=cred] (pm) at (3.0,-2.6) {PSU-M\\48 V 7.3 A\\{\tiny 기어식}};
  \draw[wire] (ac.east) -- ++(0.45,0) coordinate (t);
  \draw[wire] (t) |- (pl.west);
  \draw[wire] (t) -- (ps.west);
  \draw[wire] (t) |- (pm.west);
  % +24V_L
  \draw[rail,cblue] (pl.east) -- (13.4,2.6);
  \node[lbl,text=cblue,anchor=south west] at (4.3,2.65) {\code{+24V\_L}};
  \node[ld,anchor=south] (l1) at (11.6,3.25) {테스트 레지스터, 표시 릴레이\\감시 장치, 센서, 팬, 조작부\\F 축 브레이크(기어식, 덮개 스위치 경유)};
  \draw[wire] (11.6,2.6) -- (l1.south);
  \node[lbl,anchor=west] at (11.65,2.93) {갈래마다 퓨즈};
  % 안전 릴레이
  \node[sr] (ke) at (7.3,1.3) {K$_\mathrm{E}$};
  \node[sr] (ki) at (8.7,1.3) {K$_\mathrm{I}$};
  \draw[wire] (7.3,2.6) -- (ke.north);
  \draw[wire] (8.7,2.6) -- (ki.north);
  \node[ld,anchor=south] (lp1) at (7.0,3.25) {비상 정지 고리\\제어 커넥터 22--23};
  \node[ld,anchor=south] (lp2) at (9.0,3.25) {인터록 고리\\24--25};
  \draw[wire,dashed] (7.3,2.6) -- (7.3,3.25);
  \draw[wire,dashed] (8.7,2.6) -- (8.7,3.25);
  % +24V_SW
  \node[km] (k1) at (5.6,0) {KM1·KM2};
  \draw[rail,corange] (ps.east) -- (k1.west);
  \draw[rail,corange] (k1.east) -- (9.2,0);
  \node[lbl,text=corange,anchor=south west] at (6.3,0.05) {\code{+24V\_SW}};
  \node[ld,anchor=west] (s1) at (9.4,0.32) {퓨즈 5 A × 6 → 계수기 유닛 보드 → 시프터 4장};
  \node[ld,anchor=west] (s2) at (9.4,-0.32) {퓨즈 → 공통 계수기; 퓨즈 → K$_\mathrm{I}$ 접점 → 클럭 순서기};
  \draw[wire] (9.2,0) |- (s1.west);
  \draw[wire] (9.2,0) |- (s2.west);
  % +48V_M
  \node[km] (k2) at (5.6,-2.6) {KM1·KM2};
  \node[km] (k3) at (7.7,-2.6) {KM3};
  \draw[rail,cred] (pm.east) -- (k2.west);
  \draw[rail,cred] (k2.east) -- (k3.west);
  \draw[rail,cred] (k3.east) -- (9.8,-2.6);
  \node[ld,anchor=west] (m1) at (9.8,-2.6) {스테퍼 드라이버 (구동부의 덮개 안)};
  \node[lbl,text=cred,anchor=south west] at (8.6,-2.55) {\code{+48V\_M}};
  \node[lbl,anchor=north west,align=left,text=cgray] at (4.5,-3.05) {기어식의 \code{+24V\_SW}(1 A 이하)는 \code{+24V\_L}에서\\KM1·KM2의 다른 극을 거쳐 갈라 낸다};
  % 안전 릴레이가 접촉기를 구동
  \draw[->,cred,dashed,line width=0.6pt] (ke.south) -- (7.3,-1.3) -- (5.6,-1.3) -- (k2.north);
  \draw[->,cred,dashed,line width=0.6pt] (5.6,-1.3) -- (k1.south);
  \draw[->,cred,dashed,line width=0.6pt] (ki.south) -- (8.7,-1.7) -- (7.7,-1.7) -- (k3.north);
  \node[lbl,text=cred,anchor=north west] at (5.7,-1.36) {비상 정지: 접촉기 둘을 직렬로};
  \node[lbl,text=cred,anchor=west,align=left] at (8.8,-1.5) {덮개: 모터만 (릴레이식: 카트리지 고리, 순서기만)};
  % 0 V와 PE
  \draw[rail,cgray] (1.6,-4.4) -- (7.0,-4.4);
  \node[lbl,anchor=south west] at (1.6,-4.35) {0 V 모선: 모든 전원 장치의 0 V};
  \draw[rail,cgreen] (8.6,-4.4) -- (13.4,-4.4);
  \node[lbl,anchor=south west,text=cgreen] at (8.6,-4.35) {PE 모선: 프레임, 덮개, 케이스, 모터};
  \draw[line width=1pt,double] (7.0,-4.4) -- (8.6,-4.4);
  \node[lbl,anchor=north] at (7.8,-4.5) {연결 고리 (한 곳뿐)};
  \draw[wire,cgreen] (ac.south) |- (1.0,-4.9) -| (13.2,-4.4);
  \node[lbl,text=cgreen,anchor=south west] at (1.0,-4.88) {PE};
\end{tikzpicture}
